% !TEX program = xelatex
\documentclass{timcw}

\usepackage{amsmath, amssymb}
\usepackage{booktabs}
\usepackage{tabularx}
\usepackage{array}
\usepackage{longtable}
\usepackage{float}
\usepackage{xcolor}
\usepackage{caption}
\usepackage{tikz}
\usetikzlibrary{arrows.meta, positioning, shapes.multipart, shapes.geometric, fit, backgrounds, calc}
\usepackage{xurl}
\usepackage[hidelinks, unicode]{hyperref}

\captionsetup{labelsep=endash, font=small}

\setlength{\cftbeforechapskip}{4pt}
\setlength{\cftbeforesecskip}{1pt}
\makeatletter
\renewcommand{\cwtitle}{%
  \clearpage
  \begin{singlespacing}
  \thispagestyle{empty}

  \noindent
  \begin{minipage}[c]{0.12\textwidth}
    \includegraphics[width=0.96\linewidth]{tim-logo.png}%
  \end{minipage}%
  \hfill
  \begin{minipage}[c]{0.88\textwidth}
    \scriptsize
    \centering

    \textbf{МИНИСТЕРСТВО СЕЛЬСКОГО ХОЗЯЙСТВА
    РОССИЙСКОЙ ФЕДЕРАЦИИ}\\[0.15cm]

    ФЕДЕРАЛЬНОЕ ГОСУДАРСТВЕННОЕ БЮДЖЕТНОЕ
    ОБРАЗОВАТЕЛЬНОЕ УЧРЕЖДЕНИЕ
    ВЫСШЕГО ОБРАЗОВАНИЯ\\[0.15cm]

    \textbf{«РОССИЙСКИЙ ГОСУДАРСТВЕННЫЙ
    АГРАРНЫЙ УНИВЕРСИТЕТ --
    МСХА имени К.\,А.~ТИМИРЯЗЕВА»}\\[0.15cm]

    \textbf{(ФГБОУ ВО РГАУ --
    МСХА имени К.\,А.~Тимирязева)}
    \par
  \end{minipage}

  \par
  \medskip

  \noindent\rule{\textwidth}{0.4pt}\\[-2.0ex]
  \rule{\textwidth}{0.4pt}
  \par

  {\small
    \begin{center}
      Институт экономики и управления\\[0.1cm]
      Кафедра статистики и кибернетики\\[0.35cm]

      Направление: \cw@direction\\
      Направленность: \cw@profile\\[0.45cm]
    \end{center}
  }

  {\Large
    \begin{center}
      \textbf{КУРСОВОЙ ПРОЕКТ}
    \end{center}
  }

  {\centering
    на тему «\cw@topic».\par
  }
  \vspace{0.55cm}

  {\small
    \noindent
    \hfill
    \begin{minipage}[t]{0.55\textwidth}
      Выполнил:\\
      Студент \cw@course \ курса группы \cw@group\\
      \cw@student\\[0.8ex]

      Дата регистрации КП\\[0.6ex]
      на кафедре: \rule{4cm}{0.4pt}\\[0.15cm]

      Допущен к защите\\[0.15cm]
      Руководитель:\\[0.25cm]

      \rule{\linewidth}{0.4pt}\\
      {\footnotesize
        учёная степень, учёное звание, ФИО
      }\\[0.35cm]

      Члены комиссии:\\[0.15cm]
      \rule{0.73\linewidth}{0.4pt}\hfill
      \rule{0.23\linewidth}{0.4pt}\\[-2pt]
      {\footnotesize
        учёная степень, учёное звание, ФИО
        \hfill подпись
      }\\[0.15cm]
      \rule{0.73\linewidth}{0.4pt}\hfill
      \rule{0.23\linewidth}{0.4pt}\\[-2pt]
      {\footnotesize
        учёная степень, учёное звание, ФИО
        \hfill подпись
      }\\[0.15cm]
      \rule{0.73\linewidth}{0.4pt}\hfill
      \rule{0.23\linewidth}{0.4pt}\\[-2pt]
      {\footnotesize
        учёная степень, учёное звание, ФИО
        \hfill подпись
      }\\[0.35cm]

      Оценка \hrulefill\\[0.25cm]
      Дата защиты \hrulefill
    \end{minipage}
    \par
  }

  \vfill

  \begin{center}
    \cw@cityyear
  \end{center}

  \clearpage
  \end{singlespacing}
}
\makeatother

\cwtopic{Разработка кроссплатформенного программного средства мониторинга и оценки эффективности использования рабочего времени на предприятиях АПК}
\cwcourse{2}
\cwgroup{ДЭ 21-25}
\cwstudent{Аджемов Роман Андреевич}
\cwcityyear{Москва 2026}
\cwdirection{09.03.02 Информационные системы и технологии}
\cwprofile{Большие данные и машинное обучение}

\begin{document}
\cwtitle

\begin{singlespacing}
\tableofcontents
\end{singlespacing}
\clearpage

% =========================================================================
% ВВЕДЕНИЕ
% =========================================================================
\intro

Цифровая трансформация агропромышленного комплекса (АПК) Российской Федерации сопровождается стремительным ростом доли автоматизированных информационных процессов в ежедневной работе агрономов-аналитиков, экономистов, бухгалтеров, диспетчеров машинно-тракторного парка и специалистов по работе с государственными реестрами. Современное автоматизированное рабочее место (АРМ) сотрудника агропредприятия объединяет учётные ERP-комплексы («1С:Предприятие», «1С:Агрокомплекс»), геоинформационные системы и сервисы точного земледелия (QGIS, «АгроСигнал», Cropio, ExactFarming), а также обязательные федеральные государственные информационные системы прослеживаемости сельскохозяйственной продукции (ФГИС «Зерно», ФГИС «Сатурн», ФГИС «ВетИС.Меркурий», ЕФИС ЗСН) \cite{fgis_zerno, fgis_saturn}.

В периоды сезонных пиков сельскохозяйственных работ (посевная кампания, защита растений, уборка урожая и отгрузка партий зерна) интенсивность переключений между разнородными программными продуктами и веб-порталами резко возрастает. Это приводит к фрагментации внимания специалистов, скрытым потерям рабочего времени на рутинных операциях и снижению общей эффективности управленческих и диспетчерских подразделений. Для объективного аудита структуры трудозатрат, выявления узких мест в цифровых процессах и повышения личной продуктивности сотрудников востребованы средства автоматизированного учёта пользовательской активности (класса \emph{time-tracking}).

Вместе с тем внедрение существующих программных средств мониторинга на предприятиях АПК сталкивается с двумя системными ограничениями:
\begin{enumerate}
  \item \textbf{требования информационной безопасности и автономности.} Подавляющее большинство коммерческих облачных трекеров (RescueTime, Clockify, Hubstaff) передают заголовки открытых окон, наименования документов и историю переключений на зарубежные серверы. Для агрохолдингов и сельхозпредприятий это создаёт недопустимые риски утечки коммерчески чувствительной информации (сведений об урожайности, закупочных ценах, контрагентах, картах полей и отчётности ФГИС), а также нарушает требования к работе в изолированных контурах и на удалённых отделениях с нестабильным интернет-каналом;
  \item \textbf{гетерогенность парка операционных систем в условиях импортозамещения ПО.} Многие предприятия АПК эксплуатируют смешанную ИТ-инфраструктуру, совмещающую рабочие станции под управлением ОС Windows и современные дистрибутивы на базе ядра Linux с графическим протоколом Wayland. Большинство открытых локальных трекеров до сих пор опираются на устаревший сервер X11 либо требуют запуска громоздкого веб-интерфейса в браузере, что затрудняет их развёртывание на Linux-рабочих местах.
\end{enumerate}

\textbf{Актуальность темы} курсового проекта обусловлена необходимостью создания автономного кроссплатформенного программного средства с локальным хранением данных, адаптированного под программный профиль специалистов предприятий АПК и обеспечивающего единый функционал в ОС Windows и Linux (включая композиторы Wayland) при минимальных накладных расходах.

\textbf{Целью курсового проекта} является проектирование, программная реализация и экспериментальное исследование кроссплатформенного настольного приложения для автоматизированного учёта, классификации и визуального анализа пользовательской активности специалистов предприятий АПК.

Для достижения поставленной цели сформулированы следующие \textbf{задачи}:
\begin{enumerate}
  \item провести анализ структуры цифрового рабочего места специалиста АПК и выполнить сравнительный обзор существующих программных аналогов учёта экранного времени;
  \item исследовать низкоуровневые механизмы получения сведений об активном окне и времени бездействия пользователя (AFK) в ОС Windows (WinAPI), Linux с сервером X11 (EWMH, \texttt{libXss}) и Linux с протоколом Wayland (KWin Scripting / D-Bus, IPC композиторов Hyprland и Sway);
  \item обосновать выбор технологического стека (Python, PySide6, Qt Quick/QML, SQLite WAL) и спроектировать многослойную кроссплатформенную архитектуру приложения на основе паттерна MVVM;
  \item разработать реляционную модель базы данных, алгоритм потокового уплотнения интервалов (heartbeat-склейки) и приоритетный движок категоризации с предметным набором правил для отраслевого ПО АПК;
  \item реализовать фоновый сервис мониторинга в выделенном потоке \texttt{QThread}, слой моделей представления (\texttt{ViewModel}) и декларативный интерфейс на QML с аналитическим дашбордом, 24-часовым таймлайном, системным треем и экспортом отчётов в CSV и JSON;
  \item провести модульное, функциональное и кроссплатформенное нагрузочное тестирование на реальных стендах Windows и Linux (Wayland/KWin) с моделированием рабочего профиля специалиста АПК.
\end{enumerate}

\textbf{Объектом исследования} выступают процессы автоматизированного сбора, локального хранения и аналитической обработки данных о взаимодействии специалистов предприятий АПК с прикладным программным обеспечением на рабочих станциях.

\textbf{Предметом исследования} являются кроссплатформенные механизмы мониторинга графических оболочек ОС Windows и Linux, алгоритмы сжатия и приоритетной классификации потоков событий активности, а также архитектура настольных приложений на базе Qt~6 (PySide6/QML).

\textbf{Практическая значимость} работы заключается в создании готового к эксплуатации автономного программного инструмента \emph{Activity Tracker}, позволяющего руководителям аналитических служб и специалистам предприятий АПК объективно оценивать распределение рабочего времени между ERP, ГИС, порталами ФГИС и документооборотом без передачи данных во внешние облачные сервисы.

% =========================================================================
% ГЛАВА 1. ТЕОРЕТИЧЕСКАЯ ЧАСТЬ
% =========================================================================
\newpage
\chapter{Теоретические основы и анализ предметной области}

\section{Специфика цифрового рабочего места в АПК и обзор аналогов}

Информационная среда современного агропредприятия отличается высокой разнородностью используемых программных классов. В отличие от узкоспециализированных ИТ-команд, где значительная часть времени проходит внутри одного редактора кода, специалист информационно-аналитической, экономической или диспетчерской службы АПК в течение смены непрерывно работает в нескольких независимых контурах:
\begin{itemize}
  \item \textbf{учётные и экономические системы (ERP и табличные процессоры):} «1С:Предприятие» / «1С:Агрокомплекс», Microsoft Excel, LibreOffice Calc~--- ведение путевых листов, учёт ГСМ, сводки урожайности и себестоимости;
  \item \textbf{ГИС и платформы точного земледелия:} настольные комплексы QGIS \cite{qgis_docs} и NextGIS, а также веб-системы спутникового и телеметрического мониторинга («АгроСигнал», Cropio, ExactFarming, «История поля») для анализа индексов вегетации NDVI и выработки сельхозтехники;
  \item \textbf{федеральные государственные информационные системы АПК:} веб-порталы ФГИС «Зерно» (оформление СДИЗ) \cite{fgis_zerno}, ФГИС «Сатурн» (учёт пестицидов и агрохимикатов) \cite{fgis_saturn}, ФГИС «ВетИС.Меркурий», ЕФИС ЗСН;
  \item \textbf{электронный документооборот и коммуникации:} текстовые редакторы (Microsoft Word, LibreOffice Writer), системы ЭДО («Диадок», «СБИС»), почтовые клиенты (Thunderbird, Outlook) и корпоративные мессенджеры.
\end{itemize}

Принципиальная особенность такого профиля работы состоит в том, что существенная часть профильной нагрузки (порталы ФГИС, системы телеметрии полей, ЭДО) и одновременно потенциальных отвлечений (социальные сети, видеохостинги) сосредоточена \emph{внутри одного и того же процесса веб-браузера}. Следовательно, система учёта активности для АПК обязана классифицировать события не только по имени исполняемого файла, но и по контекстному анализу заголовка активного окна с настраиваемыми приоритетами.

Программные средства учёта рабочего времени подразделяются на три класса: системы ручного хронометража (таймеры задач), узкоспециализированные плагины телеметрии (например, для IDE) и системные автоматические трекеры активности, фиксирующие смену активного окна и периоды простоя (AFK~--- \emph{Away From Keyboard}). Разрабатываемое приложение относится к третьему классу и реализует концепцию \emph{privacy-first} (полная локальность хранения данных). В таблице~\ref{tab:analogues} приведён сравнительный анализ распространённых аналогов с учётом требований предприятий АПК.

\begin{table}[H]
\centering
\caption{Сравнительный анализ существующих систем учёта пользовательской активности}
\label{tab:analogues}
\footnotesize
\begin{tabularx}{\textwidth}{>{\raggedright\arraybackslash}p{2.5cm} >{\centering\arraybackslash}p{1.9cm} >{\centering\arraybackslash}p{2.3cm} >{\raggedright\arraybackslash}X >{\raggedright\arraybackslash}X}
\toprule
\textbf{Решение} & \textbf{Хранение данных} & \textbf{Поддержка ОС} & \textbf{Архитектура и ключевые функции} & \textbf{Ограничения для внедрения в АПК} \\
\midrule
RescueTime \cite{rescuetime} & Облако & Windows, macOS, Linux & Фоновый агент + облачный веб-портал, автокатегоризация, уведомления & Передача заголовков окон (ФГИС, договоров) на зарубежные серверы; платная подписка \\
\addlinespace
ManicTime \cite{manictime} & Локально / сервер & Windows, macOS, Linux & Нативный клиент (.NET), многодорожечный таймлайн, теги & Закрытый код; платные корпоративные лицензии; ограниченная функциональность в Linux \\
\addlinespace
ActivityWatch \cite{activitywatch} & Локально (SQLite) & Windows, macOS, Linux & Модульные процессы-наблюдатели + локальный HTTP-сервер + веб-UI & Требует браузера для работы с отчётами; под Wayland нужны разрозненные сторонние патчи \\
\addlinespace
WakaTime \cite{wakatime} & Облако & Плагины IDE & Метрики по проектам и языкам программирования & Неприменимо в АПК: не отслеживает 1С, QGIS, офисные пакеты и веб-порталы ФГИС \\
\addlinespace
\textbf{Activity Tracker (проект)} & \textbf{Локально (SQLite WAL)} & \textbf{Windows, Linux (X11 и Wayland)} & \textbf{Нативный UI (PySide6 + QML), MVVM, KWin/IPC мост, пресеты АПК, экспорт CSV/JSON} & \textbf{Отсутствие облачной синхронизации (осознанное ограничение ради приватности)} \\
\bottomrule
\end{tabularx}
\end{table}

Сопоставление систем по ключевым критериям пригодности для предприятий АПК сведено в таблице~\ref{tab:criteria_matrix}.

\begin{table}[H]
\centering
\caption{Матрица соответствия аналогов требованиям рабочих мест предприятий АПК}
\label{tab:criteria_matrix}
\footnotesize
\begin{tabularx}{\textwidth}{>{\raggedright\arraybackslash}X >{\centering\arraybackslash}p{1.8cm} >{\centering\arraybackslash}p{1.8cm} >{\centering\arraybackslash}p{2.0cm} >{\centering\arraybackslash}p{1.8cm} >{\centering\arraybackslash}p{2.1cm}}
\toprule
\textbf{Критерий} & \textbf{RescueTime} & \textbf{ManicTime} & \textbf{ActivityWatch} & \textbf{WakaTime} & \textbf{Activity Tracker} \\
\midrule
Автономная работа без внешних облаков & Нет & Частично & Да & Нет & \textbf{Да} \\
Открытый код и импортонезависимость & Нет & Нет & Да & Частично & \textbf{Да} \\
Нативный десктопный GUI без браузера & Нет & Да & Нет & Нет & \textbf{Да} \\
Поддержка Linux Wayland (KDE/Hyprland/Sway) & Частично & Частично & Частично & Н/П & \textbf{Да} \\
Разделение вкладок ФГИС / ГИС / соцсетей & Частично & Да & Да & Нет & \textbf{Да} \\
Предустановленный профиль правил АПК & Нет & Нет & Нет & Нет & \textbf{Да} \\
\bottomrule
\end{tabularx}
\end{table}

Таким образом, разработка собственного кроссплатформенного приложения с нативным интерфейсом, локальной базой данных, встроенной поддержкой Wayland и предустановленными правилами категоризации отраслевого ПО АПК является практически обоснованной задачей.

\section{Низкоуровневые механизмы мониторинга активности в Windows и Linux}

Для автоматизированного учёта экранного времени фоновая подсистема должна с шагом $\tau$ извлекать две группы системных параметров:
\begin{enumerate}
  \item \textbf{контекст активного окна}: имя исполняемого процесса (\texttt{app\_name}), путь к бинарному файлу (\texttt{executable\_path}), идентификатор процесса (\texttt{pid}) и заголовок окна (\texttt{window\_title});
  \item \textbf{длительность простоя (idle time)}: интервал в секундах с момента последнего события ввода с клавиатуры или манипулятора «мышь».
\end{enumerate}

Сложность кроссплатформенной реализации обусловлена тем, что оконные подсистемы Windows, Linux X11 и Linux Wayland основаны на принципиально различных программных интерфейсах.

\textbf{Операционная система Windows (Win32 API).} Оконный менеджер Windows предоставляет синхронный набор функций библиотеки \texttt{user32.dll} \cite{winapi_fg}:
\begin{itemize}
  \item \texttt{GetForegroundWindow()} возвращает дескриптор \texttt{HWND} окна верхнего уровня, находящегося в фокусе ввода;
  \item \texttt{GetWindowTextLengthW(hwnd)} и \texttt{GetWindowTextW(hwnd, buf, len)} считывают заголовок окна в кодировке UTF-16LE, обеспечивая корректное извлечение кириллических названий документов и порталов ФГИС;
  \item \texttt{GetWindowThreadProcessId(hwnd, \&pid)} определяет PID процесса-владельца, по которому через библиотеку \texttt{psutil} \cite{psutil_docs} извлекаются имя процесса и путь к исполняемому файлу (с обработкой исключения \texttt{AccessDenied} для окон с правами администратора);
  \item \texttt{GetLastInputInfo(\&lii)} возвращает 32-битную метку последнего события ввода \texttt{dwTime} \cite{winapi_idle}. Время простоя вычисляется относительно счётчика \texttt{kernel32.GetTickCount()} в беззнаковой арифметике по модулю $2^{32}$:
  \begin{equation}\label{eq:win_idle}
    t_{\text{idle}} = \frac{\bigl(\mathrm{GetTickCount}() - \mathrm{dwTime}\bigr) \bmod 2^{32}}{1000}.
  \end{equation}
  Маскирование по модулю $2^{32}$ (\texttt{\& 0xFFFFFFFF}) исключает переполнение и отрицательные значения при непрерывной работе рабочей станции свыше $49{,}7$ суток.
\end{itemize}

\textbf{Операционная система Linux с сервером X11.} В клиент-серверной архитектуре X Window System оконные менеджеры реализуют спецификацию EWMH (\emph{Extended Window Manager Hints}) \cite{ewmh}. Чтение атомов корневого и активного окна через \texttt{python-xlib} (\texttt{\_NET\_ACTIVE\_WINDOW}, \texttt{\_NET\_WM\_NAME} в формате \texttt{UTF8\_STRING}, \texttt{\_NET\_WM\_PID}, а также резервных \texttt{WM\_NAME} и \texttt{WM\_CLASS}) позволяет получить метаданные фокуса. Время бездействия считывается через расширение MIT-SCREEN-SAVER (\texttt{libXss.so}) вызовом \texttt{XScreenSaverQueryInfo()} \cite{xlib_manual}.

\textbf{Операционная система Linux с протоколом Wayland.} В рамках перехода на современные дистрибутивы Linux (включая среды на базе KDE Plasma~6) графический стек использует протокол Wayland, где дисплейный сервер и оконный менеджер объединены в композитор \cite{wayland}. Модель безопасности Wayland изолирует клиенты и запрещает чтение заголовков чужих окон средствами базового протокола. В связи с этим сбор данных реализуется через штатные IPC-интерфейсы композиторов:
\begin{itemize}
  \item в \textbf{KDE Plasma (KWin~5/6)} задействуется механизм сценариев \texttt{org.kde.kwin.Scripting} на сессионной шине D-Bus \cite{kwin_scripting}: временный JS-скрипт внутри KWin подписывается на активацию окна и передаёт данные в D-Bus-слот трекера;
  \item в тайловых композиторах \textbf{Hyprland} и \textbf{Sway} применяются JSON-интерфейсы \texttt{hyprctl activewindow -j} и \texttt{swaymsg -t get\_tree};
  \item определение простоя выполняется через стандартный D-Bus метод \texttt{org.freedesktop.ScreenSaver.GetSessionIdleTime()} \cite{dbus_spec} с резервным обращением к \texttt{org.gnome.Mutter.IdleMonitor.GetIdletime()}.
\end{itemize}

Сводка платформенных механизмов приведена в таблице~\ref{tab:os_mechanisms}.

\begin{table}[H]
\centering
\caption{Сопоставление механизмов сбора активности в поддерживаемых средах ОС}
\label{tab:os_mechanisms}
\footnotesize
\begin{tabularx}{\textwidth}{>{\raggedright\arraybackslash}p{2.8cm} >{\raggedright\arraybackslash}X >{\raggedright\arraybackslash}X >{\centering\arraybackslash}p{2.2cm}}
\toprule
\textbf{Среда} & \textbf{Получение активного окна} & \textbf{Определение AFK} & \textbf{Модель сбора} \\
\midrule
Windows 10/11 & \texttt{user32.GetForegroundWindow}, \texttt{GetWindowTextW}, \texttt{psutil} & \texttt{GetLastInputInfo} + \texttt{GetTickCount} & Опрос (polling) \\
\addlinespace
Linux (X11) & \texttt{python-xlib}: \texttt{\_NET\_ACTIVE\_WINDOW}, \texttt{\_NET\_WM\_NAME}, \texttt{\_NET\_WM\_PID} & \texttt{libXss}: \texttt{XScreenSaverQueryInfo} & Опрос (polling) \\
\addlinespace
Linux (Wayland / KWin) & Временный KWin-скрипт $\rightarrow$ D-Bus слот \texttt{onWindowChanged} & D-Bus \texttt{org.freedesktop.ScreenSaver} & Событийный мост + кэш \\
\addlinespace
Linux (Wayland / Hyprland, Sway) & IPC \texttt{hyprctl activewindow -j} / \texttt{swaymsg -t get\_tree} & D-Bus \texttt{ScreenSaver} / \texttt{Mutter.IdleMonitor} & Опрос (polling) \\
\bottomrule
\end{tabularx}
\end{table}

\section{Математические основы дискретного учёта и сжатия интервалов}

Автоматизированный мониторинг представляет собой дискретизацию кусочно-постоянной функции состояния пользователя $S(t) = (a(t), w(t), f(t))$ в моменты $t_0, t_1, \dots, t_N$ с периодом $\tau = t_k - t_{k-1}$. При шаге $\tau = 1$~с максимальная погрешность фиксации момента переключения не превышает $\pm 1$~с, а при независимых границах интервала математическое ожидание ошибки длительности стремится к нулю при дисперсии $\sigma^2 = \tau^2 / 6$.

Прямая запись каждого секундного отсчёта в базу данных порождает $28\,800$ строк за 8-часовую рабочую смену и $864\,000$ строк за 30 рабочих дней, приводя к росту файла БД свыше $110$--$130$~МБ в месяц и деградации скорости выборок. Применение потокового уплотнения смежных состояний (\emph{heartbeat merging}) \cite{activitywatch} позволяет сжимать поток записей пропорционально средней длительности удержания фокуса на одном окне $\bar{d}$:
\begin{equation}\label{eq:compression_theory}
  K_{\text{сж}} = \frac{\bar{d}}{\tau}.
\end{equation}
При характерном для офисно-аналитической работы в АПК значении $\bar{d} \approx 120$--$150$~с алгоритм обеспечивает теоретическое сокращение числа записей более чем в 100 раз без потери точности хронометража. Граничный порог разрыва $\Delta_{\max}$ (\texttt{pulse\_timeout}) гарантирует корректное закрытие интервала при уходе компьютера в спящий режим или временной приостановке мониторинга.

\section{Обоснование выбора технологического стека и архитектуры}

Выбор программного стека определялся требованиями кроссплатформенности, независимости от зарубежных коммерческих лицензий, простоты развёртывания на парке АРМ предприятий АПК и архитектурной чистоты \cite{martin_clean}.

\textbf{Python~3.11+ и системные библиотеки.} Язык Python обеспечивает высокую скорость реализации аналитической подсистемы и прямой доступ к системным вызовам WinAPI и \texttt{libXss} через встроенный модуль \texttt{ctypes} без необходимости компиляции C-расширений \cite{lutz_python, ramalho_python}. Библиотека \texttt{psutil} унифицирует получение метаданных процессов по PID в Windows и Linux \cite{psutil_docs}.

\textbf{PySide6 и Qt Quick (QML).} Использование официального биндинга Qt~6 для Python (\textbf{PySide6}) в связке с декларативным интерфейсом \textbf{QML} \cite{pyside6, qt_quick} имеет ряд преимуществ перед классическим \texttt{Qt Widgets} (таблица~\ref{tab:ui_comparison}).

\begin{table}[H]
\centering
\caption{Сравнение Qt Widgets и Qt Quick (QML) в рамках PySide6}
\label{tab:ui_comparison}
\footnotesize
\begin{tabularx}{\textwidth}{>{\raggedright\arraybackslash}p{3.5cm} >{\raggedright\arraybackslash}X >{\raggedright\arraybackslash}X}
\toprule
\textbf{Критерий} & \textbf{PySide6 + Qt Widgets} & \textbf{PySide6 + Qt Quick (QML)} \\
\midrule
Модель рендеринга & Растровая отрисовка на CPU (QPainter) & Аппаратный граф сцены (RHI: Direct3D, Vulkan, OpenGL) \\
Архитектурный паттерн & MVC / тесная связь UI и логики & Естественный MVVM через \texttt{@Property}, \texttt{@Slot}, \texttt{QAbstractListModel} \\
Стилизация и темы & Ограниченные таблицы стилей QSS & Декларативные биндинги палитры, мгновенная смена тем \\
Интерактивные графики & Требует QtCharts / Matplotlib & Встроенный \texttt{QtQuick.Canvas} и декларативные примитивы \\
Системный трей и D-Bus & \texttt{QSystemTrayIcon} + \texttt{QtDBus} & Поддерживаются при инициализации через \texttt{QApplication} \\
\bottomrule
\end{tabularx}
\end{table}

Встроенный в PySide6 модуль \texttt{QtDBus} позволяет работать с шиной D-Bus в Linux без установки внешней библиотеки \texttt{dbus-python}, а отказ от \texttt{QtCharts} в пользу собственных компонентов на \texttt{QtQuick.Canvas} при фиксированном стиле \texttt{QT\_QUICK\_CONTROLS\_STYLE=Basic} гарантирует идентичный внешний вид интерфейса на Windows и Linux.

\textbf{Локальная СУБД SQLite~3 в режиме WAL.} Встраиваемая реляционная СУБД SQLite~3 не требует администрирования отдельного сервера БД, хранит историю в одном файле в каталоге пользователя (\texttt{\textasciitilde/.activity\_tracker/activity.db}) и поддерживает ссылочную целостность и составные индексы \cite{sqlite_wal, kleppmann_ddia}. Включение режима \texttt{PRAGMA journal\_mode = WAL} совместно с \texttt{PRAGMA synchronous = NORMAL} обеспечивает одновременное неблокирующее чтение аналитических сводок интерфейсом и запись тиков фоновым потоком трекера.

% =========================================================================
% ГЛАВА 2. ПРАКТИЧЕСКАЯ ЧАСТЬ
% =========================================================================
\newpage
\chapter{Проектирование и программная реализация приложения}

\section{Модульная архитектура и структура проекта}

Приложение \emph{Activity Tracker} реализовано по четырёхслойной архитектуре, объединяющей паттерн \textbf{MVVM} на уровне представления данных и паттерн \textbf{«Адаптер + Фабрика»} для изоляции системно-зависимого кода ОС. Структурная схема взаимодействия компонентов представлена на рисунке~\ref{fig:architecture}.

\begin{figure}[H]
\centering
\resizebox{\textwidth}{!}{%
\begin{tikzpicture}[
  font=\sffamily\small,
  box/.style={draw=black!75, rounded corners=4pt, minimum width=3.4cm, minimum height=0.95cm, align=center, fill=#1, line width=0.7pt},
  layer/.style={draw=black!50, dashed, rounded corners=6pt, inner sep=9pt, fill=#1},
  arr/.style={-{Latex[length=2.4mm]}, thick, draw=black!80}
]
  % Слой UI (QML)
  \node[box=blue!10] (dash) {DashboardView\\(KPI, графики)};
  \node[box=blue!10, right=0.5cm of dash] (time) {TimelineView\\(лента 24 ч)};
  \node[box=blue!10, right=0.5cm of time] (rules) {RulesView\\(правила АПК)};
  \node[box=blue!10, right=0.5cm of rules] (sett) {SettingsView\\(экспорт, трей)};

  % Слой ViewModels
  \node[box=cyan!12, below=1.15cm of dash, xshift=1.95cm, minimum width=6.6cm] (dvm) {DashboardVM + AppsListModel (\texttt{QAbstractListModel})};
  \node[box=cyan!12, right=0.6cm of dvm, minimum width=6.6cm] (rvm) {RulesVM + SettingsVM (\texttt{QObject}, сигналы/слоты)};

  % Слой Core
  \node[box=green!10, below=1.15cm of dvm, xshift=-0.8cm, minimum width=4.4cm] (db) {ActivityDatabase (SQLite WAL)\\heartbeat-склейка, SQL};
  \node[box=green!10, right=0.55cm of db, minimum width=4.2cm] (cat) {ActivityCategorizer\\подстроки + Regex};
  \node[box=green!10, right=0.55cm of cat, minimum width=4.6cm] (trk) {TrackerService (\texttt{QThread})\\опрос, AFK, фильтрация};

  % Слой Platform Adapters
  \node[box=orange!12, below=1.15cm of cat, minimum width=8.4cm] (fac) {Factory + \texttt{BasePlatformAdapter} (единый контракт)};
  \node[box=yellow!16, below=0.95cm of fac, xshift=-4.6cm, minimum width=4.0cm] (win) {WindowsPlatformAdapter\\WinAPI + \texttt{psutil}};
  \node[box=yellow!16, below=0.95cm of fac, minimum width=4.0cm] (x11) {LinuxX11Adapter\\EWMH + \texttt{libXss}};
  \node[box=yellow!16, below=0.95cm of fac, xshift=4.6cm, minimum width=4.4cm] (way) {LinuxWaylandAdapter\\KWin D-Bus / Hyprland / Sway};

  % Рамки слоёв
  \begin{scope}[on background layer]
    \node[layer=blue!3, fit=(dash)(sett), label={[anchor=north west, font=\sffamily\footnotesize\bfseries]north west:View (Qt Quick / QML)}] {};
    \node[layer=cyan!4, fit=(dvm)(rvm), label={[anchor=north west, font=\sffamily\footnotesize\bfseries]north west:ViewModel (PySide6)}] {};
    \node[layer=green!3, fit=(db)(trk), label={[anchor=north west, font=\sffamily\footnotesize\bfseries]north west:Core / Model (Python)}] {};
    \node[layer=orange!4, fit=(fac)(win)(way), label={[anchor=north west, font=\sffamily\footnotesize\bfseries]north west:Platform Adapters (ОС)}] {};
  \end{scope}

  % Связи
  \draw[arr] (dash.south) -- (dvm.north -| dash.south);
  \draw[arr] (time.south) -- (dvm.north -| time.south);
  \draw[arr] (rules.south) -- (rvm.north -| rules.south);
  \draw[arr] (sett.south) -- (rvm.north -| sett.south);

  \draw[arr] (dvm.south) -- (db.north -| dvm.south);
  \draw[arr] (rvm.south -| cat.north) -- (cat.north);
  \draw[arr] (rvm.south -| trk.north) -- (trk.north);
  \draw[arr] (trk.west) -- (cat.east);
  \draw[arr] (cat.west) -- (db.east);
  \draw[arr] (trk.south) |- (fac.east);

  \draw[arr] (fac.south) -- (x11.north);
  \draw[arr] (fac.south) -| (win.north);
  \draw[arr] (fac.south) -| (way.north);
\end{tikzpicture}%
}
\caption{Многослойная архитектура приложения Activity Tracker}
\label{fig:architecture}
\end{figure}

Файловая организация прилагаемого программного проекта включает:
\begin{itemize}
  \item \texttt{main.py}~--- точка входа: инициализация \texttt{QApplication} с \texttt{setQuitOnLastWindowClosed(False)}, программная отрисовка иконки и меню системного трея (\texttt{QSystemTrayIcon}), регистрация ViewModel в \texttt{QQmlApplicationEngine}, запуск и корректная остановка фонового сервиса;
  \item \texttt{platform\_adapters/}~--- абстрактный контракт \texttt{base.py}, фабрика \texttt{factory.py}, реализации \texttt{windows.py}, \texttt{linux\_x11.py} и \texttt{linux\_wayland.py};
  \item \texttt{core/}~--- структуры данных \texttt{models.py}, движок классификации \texttt{categorizer.py}, слой SQLite и экспорта \texttt{database.py} (с пресетами ПО АПК), фоновый сервис \texttt{tracker\_service.py};
  \item \texttt{viewmodels/}~--- \texttt{apps\_model.py}, \texttt{dashboard\_vm.py}, \texttt{rules\_vm.py}, \texttt{settings\_vm.py};
  \item \texttt{qml/}~--- корневое окно \texttt{Main.qml}, компоненты \texttt{KpiCard.qml}, \texttt{DonutChart.qml}, \texttt{HourlyBarChart.qml} и четыре экрана в каталоге \texttt{views/};
  \item \texttt{tests/test\_core.py} и \texttt{benchmark.py}~--- модульные тесты ядра и стенд нагрузочного тестирования с профилем рабочей нагрузки АПК.
\end{itemize}

\section{Платформенные адаптеры и разбор KWin-скрипта для Wayland}

Общий интерфейс платформенного слоя описывается неизменяемым датаклассом \texttt{WindowInfo} и абстрактным классом \texttt{BasePlatformAdapter} в модуле \texttt{platform\_adapters/base.py}. Фабрика \texttt{create\_platform\_adapter()} в \texttt{platform\_adapters/factory.py} проверяет \texttt{sys.platform}, а в Linux дополнительно анализирует переменные \texttt{XDG\_SESSION\_TYPE} и \texttt{WAYLAND\_DISPLAY}, автоматически выбирая нужную реализацию адаптера.

В \texttt{WindowsPlatformAdapter} (\texttt{platform\_adapters/windows.py}) прямые вызовы \texttt{user32.dll} и \texttt{kernel32.dll} типизированы через \texttt{ctypes.wintypes}, что исключает неопределённое поведение на 64-битных версиях Windows. Для Linux под X11 в \texttt{LinuxX11Adapter} (\texttt{platform\_adapters/linux\_x11.py}) чтение EWMH-атомов дополнено освобождением C-структур \texttt{XFree()} и \texttt{XCloseDisplay()} в методе \texttt{close()}.

Наиболее сложной частью платформенного слоя является поддержка среды \textbf{KDE Plasma под Wayland} в модуле \texttt{platform\_adapters/linux\_wayland.py}. Поскольку Wayland запрещает сторонний опрос окон, в адаптере реализован событийно-ориентированный мост \textbf{KWin Scripting $\leftrightarrow$ QtDBus} (рисунок~\ref{fig:kwin_bridge}).

\begin{figure}[H]
\centering
\resizebox{\textwidth}{!}{%
\begin{tikzpicture}[
  font=\sffamily\footnotesize,
  comp/.style={draw=black!80, rounded corners=4pt, minimum width=3.5cm, minimum height=1.0cm, align=center, fill=#1, line width=0.75pt},
  arr/.style={-{Latex[length=2.3mm]}, thick, draw=black!85},
  darr/.style={-{Latex[length=2.3mm]}, thick, dashed, draw=blue!70!black}
]
  \node[comp=green!12] (adapter) {\textbf{LinuxWaylandAdapter}\\(инициализация / \texttt{close()})};
  \node[comp=orange!14, right=2.1cm of adapter] (dbus) {\textbf{Сессионная шина D-Bus}\\\texttt{org.kde.KWin /Scripting}\\и \texttt{ru.coursework.ActivityTracker}};
  \node[comp=blue!12, right=2.1cm of dbus] (kwin) {\textbf{Композитор KWin}\\(временный JS-скрипт\\внутри процесса KWin)};
  \node[comp=cyan!14, below=1.45cm of adapter] (recv) {\textbf{\_KWinDBusReceiver}\\(\texttt{QObject}, слот \texttt{onWindowChanged})};
  \node[comp=yellow!16, right=2.1cm of recv] (ps) {\textbf{\texttt{psutil.Process(pid)}}\\разрешение \texttt{app\_name}, \texttt{exe}};
  \node[comp=green!8, right=2.1cm of ps] (cache) {\textbf{Кэш \texttt{latest\_window}}\\(читается по \texttt{get\_active\_window()})};

  \draw[arr] (adapter.east) -- node[above, align=center]{\scriptsize 1. \texttt{registerService}\\ \scriptsize + \texttt{loadScript} / \texttt{run}} (dbus.west);
  \draw[arr] (dbus.east) -- node[above, align=center]{\scriptsize 2. Загрузка JS\\ \scriptsize и подписка} (kwin.west);
  \draw[darr] (kwin.south) |- node[pos=0.75, above, align=center]{\scriptsize 3. Событие \texttt{windowActivated} $\rightarrow$ \texttt{callDBus(cls, caption, pid)}} (dbus.south -| kwin.south) -- (dbus.south) |- (recv.east);
  \draw[arr] (recv.east) -- node[below]{\scriptsize 4. Обогащение по PID} (ps.west);
  \draw[arr] (ps.east) -- node[below]{\scriptsize 5. Обновление \texttt{WindowInfo}} (cache.west);
\end{tikzpicture}%
}
\caption{Архитектура событийного моста KWin Scripting $\leftrightarrow$ D-Bus под Wayland}
\label{fig:kwin_bridge}
\end{figure}

Работа моста организована в пять этапов:
\begin{enumerate}
  \item \textbf{Регистрация D-Bus приёмника:} объект \texttt{\_KWinDBusReceiver(QObject)} с аннотацией \texttt{@ClassInfo} регистрируется на сессионной шине под сервисом \texttt{ru.coursework.ActivityTracker} и путём \texttt{/Tracker} с экспортом слота \texttt{onWindowChanged(str, str, int)}.
  \item \textbf{Очистка остаточных скриптов:} перед загрузкой вызывается \texttt{unloadScript()} через интерфейс \texttt{org.kde.kwin.Scripting}, предотвращая дублирование после аварийного завершения.
  \item \textbf{Совместимость с Plasma~5 и Plasma~6:} генерируемый временный JavaScript-сценарий динамически проверяет наличие сигнала \texttt{workspace.windowActivated} (API KWin~6) и при его отсутствии использует \texttt{workspace.clientActivated} (API KWin~5). При смене фокуса функция \texttt{report(w)} считывает \texttt{resourceClass}, \texttt{caption} и \texttt{pid}, отправляя их через встроенный вызов \texttt{callDBus()} в слот приложения, а также однократно рапортует о текущем активном окне сразу после запуска.
  \item \textbf{Запуск сценария:} после \texttt{loadScript(path, name)} сценарий активируется вызовом \texttt{run()} на объекте \texttt{/Scripting/Script\{id\}} (KWin~6) с резервным вызовом \texttt{start()} для KWin~5.
  \item \textbf{Обогащение и очистка:} слот \texttt{onWindowChanged} уточняет имя исполняемого файла и путь через \texttt{psutil.Process(pid)} и обновляет кэш \texttt{latest\_window}, откуда \texttt{get\_active\_window()} забирает данные за $O(1)$. При выходе метод \texttt{close()} выгружает скрипт, освобождает имя D-Bus и удаляет временный файл.
\end{enumerate}

Для тайловых композиторов Hyprland и Sway в том же классе реализованы ветви опроса через \texttt{hyprctl activewindow -j} и рекурсивный обход дерева \texttt{swaymsg -t get\_tree}. Так как \texttt{QDBusConnection} требует активного цикла событий Qt, адаптер создаётся внутри рабочего потока \texttt{QThread} сервиса трекинга.

\section{Проектирование базы данных и профильная категоризация для АПК}

Реляционная база данных SQLite (модуль \texttt{core/database.py}) состоит из четырёх таблиц: \texttt{categories}, \texttt{rules}, \texttt{activity\_events} и \texttt{settings}. ER-диаграмма представлена на рисунке~\ref{fig:er_diagram}.

\begin{figure}[H]
\centering
\resizebox{\textwidth}{!}{%
\begin{tikzpicture}[
  font=\sffamily\footnotesize,
  entity/.style={
    rectangle split,
    rectangle split parts=2,
    draw=black!80,
    rounded corners=3pt,
    line width=0.8pt,
    align=left,
    rectangle split part fill={blue!12, white},
    minimum width=4.4cm
  },
  rel/.style={-{Latex[length=2.4mm]}, thick, draw=black!80}
]
  \node[entity] (cat) {
    \textbf{categories}
    \nodepart{second}
    \textbf{PK} \texttt{id} : INTEGER\\
    \texttt{name} : TEXT (UNIQUE)\\
    \texttt{color} : TEXT (HEX)\\
    \texttt{productivity} : INTEGER (-1, 0, 1)
  };

  \node[entity, right=1.6cm of cat] (rules) {
    \textbf{rules}
    \nodepart{second}
    \textbf{PK} \texttt{id} : INTEGER\\
    \texttt{pattern} : TEXT\\
    \texttt{match\_field} : TEXT (app / title)\\
    \texttt{is\_regex} : INTEGER (0 / 1)\\
    \texttt{priority} : INTEGER \quad (IDX)\\
    \textbf{FK} \texttt{category\_id} : INTEGER
  };

  \node[entity, below=1.4cm of cat, xshift=1.4cm, minimum width=5.6cm] (events) {
    \textbf{activity\_events}
    \nodepart{second}
    \textbf{PK} \texttt{id} : INTEGER\\
    \texttt{app\_name} : TEXT\\
    \texttt{window\_title} : TEXT\\
    \texttt{start\_ts} : REAL \quad (IDX)\\
    \texttt{end\_ts} : REAL \quad (IDX)\\
    \texttt{is\_afk} : INTEGER (0 / 1) \quad (IDX)\\
    \textbf{FK} \texttt{category\_id} : INTEGER
  };

  \node[entity, right=1.2cm of events, minimum width=3.8cm] (sett) {
    \textbf{settings}
    \nodepart{second}
    \textbf{PK} \texttt{key} : TEXT\\
    \texttt{value} : TEXT
  };

  \draw[rel] (rules.west) -- node[above, font=\sffamily\scriptsize]{$N : 1$ (CASCADE)} (cat.east);
  \draw[rel] (events.north -| cat.south) -- node[left, font=\sffamily\scriptsize]{$N : 1$ (SET NULL)} (cat.south);
\end{tikzpicture}%
}
\caption{ER-диаграмма реляционной базы данных SQLite}
\label{fig:er_diagram}
\end{figure}

\textbf{Алгоритм heartbeat-склейки и аналитические расчёты.} При поступлении очередного наблюдения $(a, w, f)$ в момент $t$ метод \texttt{record\_heartbeat()} расширяет последнюю запись $e_{k-1} = (a_{k-1}, w_{k-1}, f_{k-1}, t_{k-1}^{\text{start}}, t_{k-1}^{\text{end}})$ операцией \texttt{UPDATE end\_ts = max(end\_ts, t)}, если выполняется условие:
\begin{equation}\label{eq:heartbeat}
  (a = a_{k-1}) \;\wedge\; (w = w_{k-1}) \;\wedge\; (f = f_{k-1}) \;\wedge\; \bigl(0 \le t - t_{k-1}^{\text{end}} \le \Delta_{\max}\bigr).
\end{equation}
При нарушении условия~\eqref{eq:heartbeat} создаётся новая строка с началом:
\begin{equation}\label{eq:new_start}
  t_k^{\text{start}} = \begin{cases}
    t_{k-1}^{\text{end}}, & \text{если } t - \Delta t_{\text{poll}} < t_{k-1}^{\text{end}} \le t, \\[4pt]
    t - \Delta t_{\text{poll}}, & \text{иначе,}
  \end{cases}
\end{equation}
что исключает наложение границ при джиттере системного таймера. Для любого отчётного окна $[T_1, T_2]$ длительность события отсекается по границам периода на уровне SQL:
\begin{equation}\label{eq:clip}
  \Delta t_i = \max\!\Bigl(0,\; \min(t_i^{\text{end}}, T_2) - \max(t_i^{\text{start}}, T_1)\Bigr).
\end{equation}
События, пересекающие границу часа, дробятся в методе \texttt{\_distribute\_hourly()} для точного построения почасовой гистограммы. Индекс продуктивности $P$ специалиста рассчитывается по активному времени (без учёта AFK):
\begin{equation}\label{eq:productivity}
  P = \begin{cases}
    \dfrac{T_{\text{prod}}}{T_{\text{prod}} + T_{\text{neutral}} + T_{\text{distract}}} \cdot 100\%, & T_{\text{act}} > 0, \\[10pt]
    0\%, & T_{\text{act}} = 0.
  \end{cases}
\end{equation}

\textbf{Приоритетная категоризация и отраслевой профиль правил АПК.} В модуле \texttt{core/categorizer.py} реализован класс \texttt{ActivityCategorizer}, поддерживающий регистронезависимый поиск подстрок и предкомпилированные регулярные выражения (\texttt{re.compile}) с потокобезопасной перезагрузкой под \texttt{threading.RLock} и перехватом ошибок синтаксиса \texttt{re.error}. Блок-схема алгоритма классификации приведена на рисунке~\ref{fig:cat_algo}.

\begin{figure}[H]
\centering
\resizebox{0.92\textwidth}{!}{%
\begin{tikzpicture}[
  font=\sffamily\footnotesize,
  node distance=0.65cm and 1.1cm,
  term/.style={draw=black!80, rounded corners=12pt, minimum width=3.6cm, minimum height=0.75cm, align=center, fill=blue!10},
  proc/.style={draw=black!80, rounded corners=2pt, minimum width=4.6cm, minimum height=0.75cm, align=center, fill=green!8},
  dec/.style={draw=black!80, diamond, aspect=2.6, inner sep=2pt, align=center, fill=yellow!16},
  arr/.style={-{Latex[length=2.2mm]}, thick, draw=black!85}
]
  \node[term] (start) {Вход: \texttt{app\_name}, \texttt{window\_title}};
  \node[proc, below=of start] (snap) {Снять потокобезопасный снимок \texttt{\_compiled}\\под \texttt{RLock}; подготовить \texttt{app\_lower}, \texttt{title\_lower}};
  \node[dec, below=of snap] (loop) {Есть следующее правило $r_i$?};
  \node[proc, below=of loop] (field) {Выбрать целевую строку по \texttt{match\_field}\\(\texttt{app\_name} или \texttt{window\_title})};
  \node[dec, below=of field] (isre) {\texttt{r\_i.regex != None}?};
  \node[proc, left=1.2cm of isre] (rematch) {Проверка регулярным выражением\\\texttt{r\_i.regex.search(target)}};
  \node[proc, right=1.2cm of isre] (submatch) {Проверка вхождения подстроки\\\texttt{r\_i.needle\_lower in target\_lower}};
  \node[dec, below=1.15cm of isre] (hit) {Совпадение найдено?};
  \node[term, fill=green!16, below=of hit] (retcat) {Вернуть \texttt{r\_i.category\_id}};
  \node[term, fill=orange!14, right=3.8cm of loop] (retnone) {Вернуть \texttt{None}\\(«Без категории»)};

  \draw[arr] (start) -- (snap);
  \draw[arr] (snap) -- (loop);
  \draw[arr] (loop) -- node[right]{\scriptsize Да} (field);
  \draw[arr] (loop) -- node[above]{\scriptsize Нет} (retnone);
  \draw[arr] (field) -- (isre);
  \draw[arr] (isre) -- node[above]{\scriptsize Да} (rematch);
  \draw[arr] (isre) -- node[above]{\scriptsize Нет} (submatch);
  \draw[arr] (rematch.south) |- (hit.west);
  \draw[arr] (submatch.south) |- (hit.east);
  \draw[arr] (hit) -- node[right]{\scriptsize Да} (retcat);
  \draw[arr] (hit.west) -- ++(-4.3,0) |- node[pos=0.25, left]{\scriptsize Нет (следующее $r_{i+1}$)} (loop.west);
\end{tikzpicture}%
}
\caption{Блок-схема алгоритма приоритетной классификации событий}
\label{fig:cat_algo}
\end{figure}

При первичной инициализации базы данных автоматически формируется отраслевой справочник категорий и правил, охватывающий исполняемые файлы Windows и Linux, а также заголовки веб-порталов АПК (таблица~\ref{tab:apk_presets}). Выделение правил с более высоким приоритетом по полю \texttt{window\_title} позволяет однозначно различать работу в ФГИС «Зерно» или «АгроСигнал» и просмотр развлекательных сайтов внутри одного и того же браузера.

\begin{table}[H]
\centering
\caption{Базовый набор категорий и правил классификации для рабочих мест АПК}
\label{tab:apk_presets}
\footnotesize
\begin{tabularx}{\textwidth}{>{\raggedright\arraybackslash}p{3.4cm} >{\centering\arraybackslash}p{1.6cm} >{\raggedright\arraybackslash}p{2.3cm} >{\raggedright\arraybackslash}X}
\toprule
\textbf{Категория} & \textbf{Вес $p_c$} & \textbf{Поле анализа} & \textbf{Примеры шаблонов (Windows / Linux / веб)} \\
\midrule
Учёт и ERP & $+1$ & \texttt{app\_name}, \texttt{window\_title} & \texttt{1cv8}, \texttt{1cv8c}, \texttt{excel}, \texttt{scalc}; «1С:Агрокомплекс», «сводка», «ведомость», «урожайность» \\
\addlinespace
ГИС и точное земледелие & $+1$ & \texttt{app\_name}, \texttt{window\_title} & \texttt{qgis-bin}, \texttt{arcgispro}, \texttt{nextgis}; «АгроСигнал», Cropio, ExactFarming, NDVI, «карта полей» \\
\addlinespace
ФГИС и госсистемы АПК & $+1$ & \texttt{window\_title} & ФГИС «Зерно», ФГИС «Сатурн», «Меркурий», «ВетИС», «ЕФИС ЗСН», «Минсельхоз» \\
\addlinespace
Документооборот и связь & $0$ & \texttt{app\_name}, \texttt{window\_title} & \texttt{winword}, \texttt{swriter}, \texttt{thunderbird}, \texttt{telegram-desktop}; «Диадок», «СБИС», \texttt{.docx/.pdf/.ods} \\
\addlinespace
Отвлечения & $-1$ & \texttt{window\_title}, \texttt{app\_name} & YouTube, «ВКонтакте», Twitch, RuTube, Pikabu; \texttt{steam}, \texttt{discord} \\
\bottomrule
\end{tabularx}
\end{table}

При корректировке правил пользователем метод \texttt{recategorize\_all()} пересчитывает ранее накопленные интервалы единой пакетной операцией \texttt{executemany()} внутри транзакции SQLite.

\section{Фоновый сервис, диаграмма последовательности MVVM и интерфейс QML}

Фоновый мониторинг вынесен в выделенный поток \texttt{QThread} внутри фасада \texttt{TrackerService} (\texttt{core/tracker\_service.py}), исключая блокировки графического потока QML. Для предотвращения самозацикливания статистики в рабочем объекте реализован фильтр \texttt{\_is\_ignored\_window()}, исключающий собственное окно приложения как по совпадению идентификатора процесса \texttt{os.getpid()}, так и по связке процессов интерпретатора Python (\texttt{python.exe}, \texttt{pythonw.exe}, \texttt{python3}) с заголовком \texttt{Activity Tracker}.

Взаимодействие между слоями MVVM и рабочими потоками (контур периодического мониторинга и контур редактирования отраслевых правил) иллюстрирует диаграмма последовательности на рисунке~\ref{fig:mvvm_sequence}.

\begin{figure}[H]
\centering
\resizebox{\textwidth}{!}{%
\begin{tikzpicture}[
  font=\sffamily\scriptsize,
  actor/.style={draw=black!80, rounded corners=3pt, minimum width=2.25cm, minimum height=0.85cm, align=center, fill=#1, line width=0.7pt},
  lifeline/.style={dashed, draw=black!55, line width=0.6pt},
  msg/.style={-{Latex[length=2.1mm]}, thick, draw=black!85},
  ret/.style={-{Latex[length=2.1mm]}, thick, dashed, draw=blue!70!black},
  grp/.style={draw=black!45, dashed, rounded corners=4pt, fill=#1}
]
  % Участники (Lifelines)
  \node[actor=blue!12] (qml) {\textbf{QML View}\\(Main Thread)};
  \node[actor=cyan!14, right=0.55cm of qml] (vm) {\textbf{DashboardVM /}\\ \textbf{RulesVM}};
  \node[actor=green!12, right=0.55cm of vm] (db) {\textbf{ActivityDatabase}\\(SQLite WAL)};
  \node[actor=green!16, right=0.55cm of db] (trk) {\textbf{TrackerWorker}\\(\texttt{QThread})};
  \node[actor=yellow!16, right=0.55cm of trk] (cat) {\textbf{Activity-}\\ \textbf{Categorizer}};
  \node[actor=orange!14, right=0.55cm of cat] (os) {\textbf{Platform-}\\ \textbf{Adapter (ОС)}};

  % Линии жизни
  \foreach \n in {qml, vm, db, trk, cat, os} {
    \draw[lifeline] (\n.south) -- ++(0,-10.2);
  }

  % Фоновые рамки фаз
  \begin{scope}[on background layer]
    \node[grp=blue!3, fit={($(qml.south)+(-1.0,-0.25)$) ($(os.south)+(1.0,-5.9)$)}, label={[anchor=north west, font=\sffamily\scriptsize\bfseries]north west:Контур 1: Периодический цикл мониторинга (1 с) и обновление дашборда (5 с)}] {};
    \node[grp=green!3, fit={($(qml.south)+(-1.0,-6.25)$) ($(os.south)+(1.0,-9.85)$)}, label={[anchor=north west, font=\sffamily\scriptsize\bfseries]north west:Контур 2: Редактирование правил АПК пользователем и пересчёт истории}] {};
  \end{scope}

  % Контур 1: тик таймера
  \coordinate (y1) at ($(trk.south)+(0,-0.9)$);
  \draw[msg] (y1) -- node[above]{\texttt{get\_idle\_seconds()}, \texttt{get\_active\_window()}} (y1 -| os);

  \coordinate (y2) at ($(trk.south)+(0,-1.65)$);
  \draw[ret] (y2 -| os) -- node[above]{\texttt{WindowInfo}, \texttt{idle\_sec}} (y2);

  \coordinate (y3) at ($(trk.south)+(0,-2.4)$);
  \draw[msg] (y3) -- node[above]{\texttt{classify(app, title)}} (y3 -| cat);

  \coordinate (y4) at ($(trk.south)+(0,-3.1)$);
  \draw[ret] (y4 -| cat) -- node[above]{\texttt{category\_id}} (y4);

  \coordinate (y5) at ($(trk.south)+(0,-3.85)$);
  \draw[msg] (y5) -- node[above]{\texttt{record\_heartbeat(...)}} (y5 -| db);

  \coordinate (y6) at ($(trk.south)+(0,-4.6)$);
  \draw[ret] (y6) -- node[above]{Qt Signal \texttt{snapshotUpdated} (QueuedConnection)} (y6 -| vm);

  \coordinate (y7) at ($(vm.south)+(0,-5.35)$);
  \draw[ret] (y7) -- node[above]{\texttt{liveChanged} / \texttt{statsChanged} $\rightarrow$ QML bindings} (y7 -| qml);

  % Контур 2: пользователь добавляет правило
  \coordinate (y8) at ($(qml.south)+(0,-7.05)$);
  \draw[msg] (y8) -- node[above]{\texttt{rulesVM.addRule(...)} / \texttt{recategorizeHistory()}} (y8 -| vm);

  \coordinate (y9) at ($(vm.south)+(0,-7.8)$);
  \draw[msg] (y9) -- node[above]{\texttt{add\_rule()}, \texttt{recategorize\_all()}} (y9 -| db);

  \coordinate (y10) at ($(vm.south)+(0,-8.55)$);
  \draw[msg] (y10) -- node[above]{Signal \texttt{rulesUpdated} $\rightarrow$ \texttt{reload\_rules()}} (y10 -| trk);

  \coordinate (y11) at ($(trk.south)+(0,-9.3)$);
  \draw[msg] (y11) -- node[above]{\texttt{reload\_rules(db.get\_rules())}} (y11 -| cat);
\end{tikzpicture}%
}
\caption{Диаграмма последовательности взаимодействия слоёв MVVM и потоков}
\label{fig:mvvm_sequence}
\end{figure}

Для оптимизации нагрузки в \texttt{DashboardVM} применён двухуровневый цикл обновления: индикатор текущего окна в боковой панели обновляется на каждом тике (1~с), а агрегация графиков и модели \texttt{AppsListModel} выполняется раз в 5~с либо мгновенно при смене периода.

Графический интерфейс (\texttt{qml/}) включает четыре экрана и ряд инженерных решений для стабильной работы в Qt~6 \cite{qt_quick}:
\begin{itemize}
  \item \textbf{Дашборд (\texttt{DashboardView.qml}):} выбор периода (сегодня, вчера, 7 и 30 дней), 4 KPI-карточки, кольцевая диаграмма категорий \texttt{DonutChart} на \texttt{QtQuick.Canvas}, почасовая диаграмма \texttt{HourlyBarChart} (общая и продуктивная активность) и рейтинг приложений. В \texttt{HourlyBarChart} применён единый \texttt{MouseArea} поверх всех 24 слотов с \texttt{spacing: 0} и внутренними боковыми полями столбцов, что полностью устраняет мерцание подсказки при прохождении курсора через зазоры между часами;
  \item \textbf{Таймлайн дня (\texttt{TimelineView.qml}):} 24-часовая цветовая шкала с явным объявлением \texttt{required property int index} в делегатах меток времени (требование строгой типизации Qt~6) и подробный журнал интервалов;
  \item \textbf{Категории и правила (\texttt{RulesView.qml}):} добавление и удаление категорий и правил (подстроки/Regex), кнопка пересчёта накопленной истории;
  \item \textbf{Настройки и трей (\texttt{SettingsView.qml}, \texttt{Main.qml}):} настройка интервала опроса и порога AFK, мгновенное переключение тёмной и светлой тем, экспорт отчётов за выбранный период в форматы CSV (UTF-8-SIG с разделителем «\texttt{;}» для прямого открытия в Excel и LibreOffice Calc) и JSON, а также сворачивание окна в системный трей (\texttt{QSystemTrayIcon}) без остановки сбора данных.
\end{itemize}

% TODO: после снятия скриншотов положите файлы в папку images/ рядом с main.tex
% и замените блоки \fbox на \includegraphics.
\begin{figure}[H]
\centering
% \includegraphics[width=0.95\textwidth]{images/dashboard_screenshot.png}
\fbox{\parbox{0.9\textwidth}{\centering\vspace{2.0cm}\textit{[Место для скриншота главного экрана «Дашборд» с отраслевыми категориями АПК]}\vspace{2.0cm}}}
\caption{Главный экран приложения --- аналитический дашборд}
\label{fig:ui_dashboard}
\end{figure}

\begin{figure}[H]
\centering
% \includegraphics[width=0.95\textwidth]{images/timeline_screenshot.png}
\fbox{\parbox{0.9\textwidth}{\centering\vspace{1.8cm}\textit{[Место для скриншота экрана «Таймлайн дня» с 24-часовой цветовой шкалой]}\vspace{1.8cm}}}
\caption{Экран хронологической ленты активности за выбранный день}
\label{fig:ui_timeline}
\end{figure}

\begin{figure}[H]
\centering
% \includegraphics[width=0.95\textwidth]{images/rules_settings_screenshot.png}
\fbox{\parbox{0.9\textwidth}{\centering\vspace{1.8cm}\textit{[Место для скриншота экрана «Категории и правила» / «Настройки и экспорт»]}\vspace{1.8cm}}}
\caption{Экран управления правилами классификации ПО АПК и параметрами трекинга}
\label{fig:ui_rules}
\end{figure}

% =========================================================================
% ГЛАВА 3. ТЕСТИРОВАНИЕ И ОЦЕНКА ПРОИЗВОДИТЕЛЬНОСТИ
% =========================================================================
\newpage
\chapter{Тестирование и оценка производительности}

\section{Методика тестирования и конфигурация стендов}

Верификация разработанного программного комплекса проводилась в три этапа: модульное тестирование ядра (\texttt{tests/test\_core.py}), функциональная проверка пользовательских сценариев и кроссплатформенное нагрузочное тестирование с помощью бенчмарка \texttt{benchmark.py}, воспроизводящего профиль рабочей нагрузки специалиста предприятия АПК.

\textbf{Модульное тестирование (\texttt{unittest})} выполняется на изолированной базе SQLite в оперативной памяти (\texttt{:memory:}) и проверяет: приоритетную классификацию внутри процесса браузера, склейку смежных тиков и разрыв интервала при паузе свыше $\Delta_{\max}$, отсечение событий по формуле~\eqref{eq:clip}, расчёт индекса продуктивности~\eqref{eq:productivity} и корректность экспорта в CSV и JSON.

\textbf{Нагрузочный стенд \texttt{benchmark.py}} включает четыре воспроизводимые серии замеров:
\begin{enumerate}
  \item \textbf{живой опрос платформенного адаптера} (30~с, шаг 1~с) с фиксацией среднего времени, 95-го перцентиля (\texttt{p95}) и максимума системного вызова, загрузки CPU и объёма резидентной памяти (RSS);
  \item \textbf{эффективность heartbeat-склейки} на синтетическом потоке за 30 рабочих дней по 8~часов ($864\,000$ секундных тиков по приложениям АПК: «1С:Агрокомплекс», LibreOffice Calc/Writer, QGIS, «АгроСигнал», ФГИС «Зерно», ФГИС «Сатурн», «ВетИС.Меркурий», Thunderbird, Telegram и отвлечения) с последующим сбросом WAL-журнала (\texttt{PRAGMA wal\_checkpoint(TRUNCATE)});
  \item \textbf{латентность аналитических SQL-запросов} (50 повторений после прогрева кэша для 4 типовых выборок на 30-дневной базе);
  \item \textbf{пропускная способность категоризатора} при обработке $10\,000$ событий набором из 25 отраслевых правил (комбинация подстрок и регулярных выражений).
\end{enumerate}

Испытания проведены на двух физических рабочих станциях, соответствующих типовой гетерогенной инфраструктуре предприятия (настольный ПК под Windows и мобильная рабочая станция под Linux с Wayland). Характеристики стендов сведены в таблице~\ref{tab:testbeds}.

\begin{table}[H]
\centering
\caption{Конфигурация тестовых стендов}
\label{tab:testbeds}
\footnotesize
\begin{tabularx}{\textwidth}{>{\raggedright\arraybackslash}p{4.0cm} >{\raggedright\arraybackslash}X >{\raggedright\arraybackslash}X}
\toprule
\textbf{Параметр стенда} & \textbf{Стенд №1 (Windows)} & \textbf{Стенд №2 (Linux)} \\
\midrule
Операционная система & Windows 11 26H2 & Arch Linux (ядро 7.3.6-arch1-1) \\
Графическая подсистема & Win32 / DWM (WinAPI) & Wayland (KDE Plasma / KWin) \\
Процессор (CPU) & AMD Ryzen 5 5600X (6 ядер / 12 потоков) & AMD Ryzen 5 5500U (6 ядер / 12 потоков) \\
Оперативная память & 16~ГБ RAM & 16~ГБ RAM \\
Среда исполнения & Python 3.14.7, PySide6 6.12.0 & Python 3.14.10, PySide6 6.12.0 \\
\bottomrule
\end{tabularx}
\end{table}

\section{Анализ результатов функционального и нагрузочного тестирования}

Автоматизированные модульные тесты ядра и ручная проверка сценариев интерфейса пройдены полностью на обоих тестовых стендах (таблица~\ref{tab:func_tests}).

\begin{table}[H]
\centering
\caption{Результаты функционального тестирования приложения}
\label{tab:func_tests}
\footnotesize
\begin{tabularx}{\textwidth}{>{\centering\arraybackslash}p{0.7cm} >{\raggedright\arraybackslash}p{4.5cm} >{\raggedright\arraybackslash}X >{\centering\arraybackslash}p{2.1cm}}
\toprule
\textbf{№} & \textbf{Сценарий проверки} & \textbf{Ожидаемый результат} & \textbf{Статус} \\
\midrule
1 & Переключение окон в Windows~11 и Linux (Wayland/KWin) & Корректное извлечение имени процесса, PID и заголовка UTF-8 & Пройден \\
\addlinespace
2 & Превышение порога бездействия (AFK) & Переход в состояние \texttt{is\_afk = 1}, возврат в активный режим при вводе & Пройден \\
\addlinespace
3 & Фильтрация собственного окна & Окно «Activity Tracker» (\texttt{python.exe} / \texttt{os.getpid()}) не попадает в БД & Пройден \\
\addlinespace
4 & Разделение вкладок ФГИС, ГИС и соцсетей в браузере & Корректное разнесение вкладок «ФГИС Зерно», «АгроСигнал» и YouTube по разным категориям & Пройден \\
\addlinespace
5 & Heartbeat-склейка и отсечение границ суток & Сжатие смежных тиков в один интервал; точный расчёт при переходе через 00:00 & Пройден \\
\addlinespace
6 & Интерактивность графиков и переключение темы & Плавный hover без мерцания на \texttt{HourlyBarChart}; перерисовка \texttt{DonutChart} & Пройден \\
\addlinespace
7 & Системный трей и экспорт CSV/JSON & Непрерывный сбор в фоне; корректное открытие CSV в Excel и LibreOffice Calc & Пройден \\
\bottomrule
\end{tabularx}
\end{table}

Фактические результаты нагрузочного тестирования, полученные скриптом \texttt{benchmark.py} на обоих стендах, сведены в таблице~\ref{tab:benchmark_cmp}. Визуальное сравнение задержек выполнения аналитических SQL-запросов представлено на рисунке~\ref{fig:sql_bench_chart}.

\begin{table}[H]
\centering
\caption{Сравнительные результаты нагрузочного тестирования на Windows и Linux (профиль АПК)}
\label{tab:benchmark_cmp}
\footnotesize
\begin{tabularx}{\textwidth}{>{\raggedright\arraybackslash}X >{\centering\arraybackslash}p{3.1cm} >{\centering\arraybackslash}p{3.1cm}}
\toprule
\textbf{Измеряемый показатель} & \textbf{Windows 11 (WinAPI)} & \textbf{Linux (Wayland / KWin)} \\
\midrule
Среднее время опроса окна и AFK, мс & 0{,}160 & 1{,}119 \\
95-й перцентиль времени опроса, мс & 0{,}323 & 3{,}056 \\
Максимальное время опроса, мс & 0{,}414 & 3{,}223 \\
Средняя / пиковая загрузка CPU в фоне, \% & $<0{,}01$ / $<0{,}01$ & 0{,}14 / 1{,}00 \\
Среднее / пиковое потребление RAM (RSS), МБ & 26{,}26 / 26{,}30 & 46{,}43 / 46{,}45 \\
Число строк БД за 30 дн. (без / со склейкой) & 864\,000 / 6\,066 & 864\,000 / 6\,066 \\
Коэффициент сжатия строк (heartbeat) & 142{,}43$\times$ & 142{,}43$\times$ \\
Размер БД за 30 дн. (без / со склейкой), МБ & 113{,}27 / 0{,}84 & 127{,}70 / 0{,}94 \\
Снижение объёма файла БД, \% & 99{,}26 & 99{,}27 \\
Агрегация по категориям за 1 день (avg / p95), мс & 0{,}790 / 0{,}866 & 1{,}843 / 1{,}863 \\
Топ-10 приложений за 7 дней (avg / p95), мс & 0{,}655 / 0{,}741 & 1{,}507 / 1{,}523 \\
Сводка по категориям за 30 дней (avg / p95), мс & 1{,}360 / 1{,}592 & 3{,}115 / 3{,}156 \\
Выборка ленты таймлайна за 1 день (avg / p95), мс & 0{,}473 / 0{,}583 & 1{,}022 / 1{,}036 \\
Скорость категоризации (25 правил), соб./с & 194\,852 & 98\,343 \\
\bottomrule
\end{tabularx}
\end{table}

\begin{figure}[H]
\centering
\resizebox{0.94\textwidth}{!}{%
\begin{tikzpicture}[font=\sffamily\footnotesize]
  % Сетка и ось Y (масштаб: 1 мс = 1.25 см, максимум 3.5 мс)
  \draw[->, thick] (0,0) -- (0,4.65) node[above] {Время, мс};
  \draw[->, thick] (0,0) -- (12.8,0);

  \foreach \y/\lbl in {0/0, 0.5/0{,}5, 1.0/1{,}0, 1.5/1{,}5, 2.0/2{,}0, 2.5/2{,}5, 3.0/3{,}0, 3.5/3{,}5} {
    \draw[gray!30, dashed] (0, \y*1.25) -- (12.5, \y*1.25);
    \draw (0, \y*1.25) -- (-0.12, \y*1.25) node[left] {\lbl};
  }

  % Группа 1: Агрегация (1 день) — 0.790 vs 1.843
  \fill[blue!65!cyan!80, draw=black!70] (0.95, 0) rectangle (1.70, 0.790*1.25);
  \node[above, font=\sffamily\scriptsize] at (1.325, 0.790*1.25) {0{,}79};
  \fill[orange!80!red!75, draw=black!70] (1.80, 0) rectangle (2.55, 1.843*1.25);
  \node[above, font=\sffamily\scriptsize] at (2.175, 1.843*1.25) {1{,}84};
  \node[below, align=center] at (1.75, -0.15) {Категории\\(1 день)};

  % Группа 2: Топ-10 (7 дней) — 0.655 vs 1.507
  \fill[blue!65!cyan!80, draw=black!70] (3.95, 0) rectangle (4.70, 0.655*1.25);
  \node[above, font=\sffamily\scriptsize] at (4.325, 0.655*1.25) {0{,}66};
  \fill[orange!80!red!75, draw=black!70] (4.80, 0) rectangle (5.55, 1.507*1.25);
  \node[above, font=\sffamily\scriptsize] at (5.175, 1.507*1.25) {1{,}51};
  \node[below, align=center] at (4.75, -0.15) {Топ-10 ПО\\(7 дней)};

  % Группа 3: Сводка (30 дней) — 1.360 vs 3.115
  \fill[blue!65!cyan!80, draw=black!70] (6.95, 0) rectangle (7.70, 1.360*1.25);
  \node[above, font=\sffamily\scriptsize] at (7.325, 1.360*1.25) {1{,}36};
  \fill[orange!80!red!75, draw=black!70] (7.80, 0) rectangle (8.55, 3.115*1.25);
  \node[above, font=\sffamily\scriptsize] at (8.175, 3.115*1.25) {3{,}12};
  \node[below, align=center] at (7.75, -0.15) {Сводка\\(30 дней)};

  % Группа 4: Таймлайн (1 день) — 0.473 vs 1.022
  \fill[blue!65!cyan!80, draw=black!70] (9.95, 0) rectangle (10.70, 0.473*1.25);
  \node[above, font=\sffamily\scriptsize] at (10.325, 0.473*1.25) {0{,}47};
  \fill[orange!80!red!75, draw=black!70] (10.80, 0) rectangle (11.55, 1.022*1.25);
  \node[above, font=\sffamily\scriptsize] at (11.175, 1.022*1.25) {1{,}02};
  \node[below, align=center] at (10.75, -0.15) {Таймлайн\\(1 день)};

  % Легенда
  \fill[blue!65!cyan!80, draw=black!70] (0.95, 4.15) rectangle (1.35, 4.42);
  \node[right] at (1.42, 4.285) {Windows 11 (Ryzen 5 5600X)};
  \fill[orange!80!red!75, draw=black!70] (6.25, 4.15) rectangle (6.65, 4.42);
  \node[right] at (6.72, 4.285) {Arch Linux / Wayland (Ryzen 5 5500U)};
\end{tikzpicture}%
}
\caption{Сравнение среднего времени выполнения аналитических SQL-запросов (50 прогонов)}
\label{fig:sql_bench_chart}
\end{figure}

Детальный анализ экспериментальных данных позволяет сформулировать следующие технические выводы:
\begin{itemize}
  \item \textbf{Низкие накладные расходы фонового мониторинга.} На Windows~11 прямые вызовы WinAPI через \texttt{ctypes} выполняются в среднем за 0{,}160~мс (p95~--- 0{,}323~мс), обеспечивая загрузку CPU ниже порога погрешности измерения ($<0{,}01\%$) при потреблении 26{,}26~МБ оперативной памяти. На Linux (Wayland/KWin) среднее время цикла опроса составляет 1{,}119~мс (p95~--- 3{,}056~мс), а объём RSS~--- 46{,}43~МБ: увеличение относительно Windows объясняется загрузкой библиотек \texttt{QtCore}/\texttt{QtDBus} и межпроцессным обменом с сервисом \texttt{org.freedesktop.ScreenSaver} по шине D-Bus, а также более низкими тактовыми частотами энергоэффективного мобильного процессора Ryzen~5~5500U. При этом средняя загрузка CPU на Linux составляет всего 0{,}14\%, что не оказывает заметного влияния на автономность ноутбука и отзывчивость рабочего стола.
  \item \textbf{Высокая эффективность алгоритма heartbeat-склейки.} На 30-дневном потоке из 864\,000 секундных наблюдений алгоритм уплотнителя сокращает таблицу до 6\,066 записей (коэффициент сжатия $142{,}43\times$), полностью подтверждая теоретическую оценку~\eqref{eq:compression_theory}. Физический размер базы данных уменьшается со 113{,}27 до 0{,}84~МБ на Windows (снижение на 99{,}26\%) и со 127{,}70 до 0{,}94~МБ на Linux (снижение на 99{,}27\%). Таким образом, годовая история активности сотрудника предприятия АПК занимает порядка 10--12~МБ на диске.
  \item \textbf{Мгновенная аналитика и категоризация.} Даже наиболее ресурсоёмкий запрос~--- агрегация 30-дневной статистики по категориям~--- выполняется в среднем за 1{,}360~мс на Windows и 3{,}115~мс на мобильном процессоре Linux, что в разы меньше длительности одного кадра интерфейса (16{,}6~мс при 60~Гц и 8{,}3~мс при 120~Гц). Пропускная способность движка категоризации с 25 отраслевыми правилами АПК достигает 194\,852 соб./с на Ryzen~5~5600X и 98\,343 соб./с на Ryzen~5~5500U, позволяя переклассифицировать накопленную историю за месяц менее чем за 0{,}1~с. Двукратная разница в скорости вычислений и SQL-запросов между стендами закономерно обусловлена отличием однопоточной производительности настольного процессора архитектуры Zen~3 (5600X) и мобильного процессора с теплопакетом 15~Вт архитектуры Zen~2 (5500U).
\end{itemize}

% =========================================================================
% ЗАКЛЮЧЕНИЕ
% =========================================================================
\newpage
\conclusion

В ходе выполнения курсового проекта спроектировано, программно реализовано и экспериментально протестировано кроссплатформенное настольное приложение \emph{Activity Tracker} для автоматизированного учёта, классификации и анализа пользовательской активности специалистов предприятий агропромышленного комплекса.

В результате проделанной работы выполнены все поставленные задачи:
\begin{enumerate}
  \item проанализирована структура цифрового рабочего места специалиста АПК (ERP «1С:Агрокомплекс», ГИС-платформы точного земледелия, порталы ФГИС «Зерно» / «Сатурн» / «Меркурий», системы ЭДО) и обоснована потребность в автономном локальном решении без облачной телеметрии;
  \item исследованы и изолированы за общим программным интерфейсом платформенные механизмы сбора активности Windows (WinAPI), Linux X11 (EWMH и \texttt{libXss}) и Linux Wayland (событийный мост KWin Scripting $\leftrightarrow$ D-Bus, IPC Hyprland и Sway);
  \item реализована четырёхслойная модульная архитектура на базе паттерна MVVM, встраиваемая БД SQLite в режиме WAL, алгоритм heartbeat-склейки интервалов и приоритетный движок категоризации с предустановленными правилами для программного обеспечения АПК;
  \item разработан нативный графический интерфейс на PySide6 и Qt Quick (QML), включающий аналитический дашборд с кастомными графиками, 24-часовой таймлайн, редактор категорий и правил с пересчётом истории, системный трей, светлую/тёмную тему и экспорт отчётов в CSV и JSON;
  \item проведены модульные, функциональные и нагрузочные испытания на реальных стендах Windows~11 и Arch~Linux (Wayland/KWin), подтвердившие минимальную нагрузку на процессор ($<0{,}01\%$ на Windows и 0{,}14\% на Linux), сжатие локальной БД в 142{,}43 раза (свыше 99{,}26\% экономии дискового пространства) и время аналитических SQL-выборок в пределах 0{,}47--3{,}12~мс.
\end{enumerate}

Перспективами развития проекта являются внедрение пользовательского чёрного списка исключаемых окон в интерфейсе настроек, расширение поддержки Wayland-композиторов (GNOME Shell) и отечественных корпоративных дистрибутивов Linux (Astra Linux, Альт Линукс, РЕД ОС), а также добавление модуля обезличенной агрегированной выгрузки показателей по подразделениям агропредприятия.

% =========================================================================
% СПИСОК ИСПОЛЬЗОВАННЫХ ИСТОЧНИКОВ
% =========================================================================
\clearpage
\centeredrefname
\addcontentsline{toc}{chapter}{Список использованных источников}
\begin{thebibliography}{99}

\bibitem{fgis_zerno}
Федеральная государственная информационная система прослеживаемости зерна и продуктов переработки зерна (ФГИС «Зерно») [Электронный ресурс] / ФГБУ «Центр Агроаналитики» Минсельхоза России. --- Режим доступа: \url{https://specagro.ru/fgis} (дата обращения: 15.11.2026).

\bibitem{fgis_saturn}
Федеральная государственная информационная система прослеживаемости пестицидов и агрохимикатов (ФГИС «Сатурн») [Электронный ресурс] / Россельхознадзор. --- Режим доступа: \url{https://fgis-saturn.ru/} (дата обращения: 15.11.2026).

\bibitem{qgis_docs}
QGIS: A Free and Open Source Geographic Information System. Documentation [Электронный ресурс]. --- Режим доступа: \url{https://docs.qgis.org/} (дата обращения: 15.11.2026).

\bibitem{rescuetime}
RescueTime: Automatic Time Tracking Software [Электронный ресурс]. --- Режим доступа: \url{https://www.rescuetime.com/} (дата обращения: 15.11.2026).

\bibitem{manictime}
ManicTime: Time Tracking Software [Электронный ресурс]. --- Режим доступа: \url{https://www.manictime.com/} (дата обращения: 15.11.2026).

\bibitem{activitywatch}
ActivityWatch: Open-Source Automated Time Tracker [Электронный ресурс]. --- Режим доступа: \url{https://activitywatch.net/} (дата обращения: 15.11.2026).

\bibitem{wakatime}
WakaTime: Dashboards for Developers [Электронный ресурс]. --- Режим доступа: \url{https://wakatime.com/} (дата обращения: 15.11.2026).

\bibitem{winapi_fg}
Microsoft Learn. GetForegroundWindow function (winuser.h) [Электронный ресурс]. --- Режим доступа: \url{https://learn.microsoft.com/en-us/windows/win32/api/winuser/nf-winuser-getforegroundwindow} (дата обращения: 16.11.2026).

\bibitem{winapi_idle}
Microsoft Learn. GetLastInputInfo function (winuser.h) [Электронный ресурс]. --- Режим доступа: \url{https://learn.microsoft.com/en-us/windows/win32/api/winuser/nf-winuser-getlastinputinfo} (дата обращения: 16.11.2026).

\bibitem{ewmh}
freedesktop.org. Extended Window Manager Hints (EWMH) Specification [Электронный ресурс]. --- Режим доступа: \url{https://specifications.freedesktop.org/wm-spec/latest/} (дата обращения: 16.11.2026).

\bibitem{xlib_manual}
Gettys J., Scheifler R. W. Xlib --- C Language X Interface [Электронный ресурс]. --- Режим доступа: \url{https://www.x.org/releases/current/doc/libX11/libX11/libX11.html} (дата обращения: 16.11.2026).

\bibitem{wayland}
Wayland Architecture and Protocol Documentation [Электронный ресурс]. --- Режим доступа: \url{https://wayland.freedesktop.org/docs/html/} (дата обращения: 17.11.2026).

\bibitem{kwin_scripting}
KDE Developer Documentation. KWin Scripting API [Электронный ресурс]. --- Режим доступа: \url{https://develop.kde.org/docs/plasma/kwin/} (дата обращения: 17.11.2026).

\bibitem{dbus_spec}
freedesktop.org. D-Bus Specification [Электронный ресурс]. --- Режим доступа: \url{https://dbus.freedesktop.org/doc/dbus-specification.html} (дата обращения: 17.11.2026).

\bibitem{pyside6}
Qt for Python (PySide6) Official Documentation [Электронный ресурс]. --- Режим доступа: \url{https://doc.qt.io/qtforpython-6/} (дата обращения: 18.11.2026).

\bibitem{qt_quick}
Qt Group. Qt Quick and QML Applications Documentation [Электронный ресурс]. --- Режим доступа: \url{https://doc.qt.io/qt-6/qtquick-index.html} (дата обращения: 18.11.2026).

\bibitem{sqlite_wal}
SQLite Documentation. Write-Ahead Logging [Электронный ресурс]. --- Режим доступа: \url{https://www.sqlite.org/wal.html} (дата обращения: 18.11.2026).

\bibitem{psutil_docs}
Rodola G. psutil documentation: cross-platform process and system utilities [Электронный ресурс]. --- Режим доступа: \url{https://psutil.readthedocs.io/} (дата обращения: 18.11.2026).

\bibitem{martin_clean}
Мартин Р. Чистая архитектура. Искусство разработки программного обеспечения / Р.~Мартин ; пер. с англ. --- СПб. : Питер, 2021. --- 352~с.

\bibitem{lutz_python}
Лутц М. Изучаем Python : в 2 т. / М.~Лутц ; пер. с англ. --- 5-е изд. --- СПб. : Диалектика, 2020. --- Т.~1. --- 832~с.

\bibitem{ramalho_python}
Рамальо Л. Python. К вершинам мастерства / Л.~Рамальо ; пер. с англ. --- 2-е изд. --- М. : ДМК Пресс, 2022. --- 898~с.

\bibitem{kleppmann_ddia}
Клеппман М. Высоконагруженные приложения. Программирование, масштабирование, поддержка / М.~Клеппман ; пер. с англ. --- СПб. : Питер, 2021. --- 640~с.

\end{thebibliography}

\end{document}